\section{Plethysm coefficients: the modular index}\label{sec:plethysm}

Let $V=\C^n$ with coordinates $x_1,\dots,x_n$, $n\ge4$, $\nu=\gcd(n,3)$,
$d_0=n/\nu$, $\gamma=3/\nu$ (an odd integer) and $M=2^{n-1}$. Points of
$\Sym^3V^*$ are cubic forms, with $(a\cdot F)(x)=F(a^{-1}x)$. Let
$R=\C[\Sym^3V^*]^{SL_n}$ with its coefficient grading $R_D$.

\paragraph{The multiplicity algebra.}
If $R_D\ne0$, the scalar $\zeta I_n\in SL_n$, $\zeta^n=1$, acts on $R_D$ by
$\zeta^{\pm3D}$, so $n\mid3D$, that is $D=d_0s$. Writing $a=zh$ with
$h\in SL_n$ and $z^n=\det a$, the element $a\in GL_n$ acts on $R_{d_0s}$ by the
character $(\det a)^{\gamma s}$, and, as functions,
\begin{equation}\label{eq:detcharacter}
 P(a\cdot F)=(\det a)^{-\gamma s}P(F)\qquad(P\in R_{d_0s}).
\end{equation}
Choosing a volume vector $\varepsilon\in\Lambda^nV$, evaluation of an
intertwiner at $\varepsilon^{\otimes\gamma s}$ identifies the degree-$s$ part
of the multiplicity algebra $\bigoplus_s\Hom_{GL_n}((\det)^{\gamma s},
\Sym^{d_0s}(\Sym^3V))$ of \cite[Section 4.1.2, eq.~(4.5)]{GCT6} with $R_{d_0s}$,
compatibly with products. So the multiplicity algebra is
$T^{(n)}=\bigoplus_sR_{d_0s}$, the whole invariant ring with its degrees
divided by $d_0$, a finitely generated positively graded domain with
$T^{(n)}_0=\C$, and $H_n=\Hilb_{T^{(n)}}$. Let $X=\Spec T^{(n)}$. By
\eqref{eq:detcharacter}, $a\in GL_n$ acts on $X$ as the grading element
$(\det a)^{-\gamma}$; as $\det$ is onto $\C^*$ and $\gamma$ is odd, every
grading element arises.

\begin{proposition}\label{prop:chaincubic}
Let $F_n=x_1^2x_2+x_2^2x_3+\dots+x_{n-1}^2x_n+x_n^3$ and $p_n\in X$ its image.
\begin{enumerate}
\item $F_n$ defines a smooth hypersurface, and its $SL_n$-orbit is closed.
\item Every linear transformation preserving $F_n$ up to a scalar is a scalar
multiple of a diagonal matrix $\diag(a_1,a_1^{-2},a_1^4,\dots,a_1^{(-2)^{n-1}})$
with $a_1^{3M}=1$. These matrices form the stabilizer of $F_n$ in $GL_n$, a
cyclic group of order $3M$; the stabilizer in $SL_n$ is the scalar group
$\mu_\nu$, which acts trivially on cubic forms.
\item The stabilizer of $p_n$ in the grading group is exactly $\mu_M$.
\item $X$ is smooth at $p_n$, and for a primitive $M$-th root of unity $\zeta$
the fixed scheme $X^\zeta$ is smooth of dimension one at $p_n$.
\end{enumerate}
\end{proposition}

\begin{proof}
(1) The partial derivatives are $2x_1x_2$, $x_{i-1}^2+2x_ix_{i+1}$
($1<i<n$) and $x_{n-1}^2+3x_n^2$. At a common zero $x\ne0$, let $j$ be least
with $x_j\ne0$. If $j<n$, the $j$-th derivative gives $x_{j+1}=0$, and then
the $(j+1)$-st is $x_j^2\ne0$; if $j=n$, the last one is $3x_n^2\ne0$. So
$F_n$ is smooth. For the closed orbit, put $r_i=(1-(-1/2)^i)/3$ for $i<n$ and
$r_n=(1-r_{n-1})/3$; all are positive, and with $a_i=2e_i+e_{i+1}$ ($i<n$),
$a_n=3e_n$, one has $\sum_ir_ia_i=\one$. Fix an $SU_n$-invariant Hermitian
product on $\Sym^3V^*$; distinct monomials are orthogonal, having distinct
torus weights. The cubic $F=\sum_i(r_i/\|m_i\|^2)^{1/2}m_i$, where $m_i$ are the
chain monomials, is a zero of the moment map of $SU_n$: its diagonal part is the restriction of $\sum_ir_ia_i$ to the trace-free
diagonal matrices, which is zero because $\sum_ir_ia_i=\one$, and its off-diagonal part
vanishes because no two of the weights $a_i$ differ by a root $e_a-e_b$
(their differences have an entry $\pm2$, or the form $2e_i-e_{i+1}-e_{i+2}$, or
$2e_{n-1}-2e_n$). By the Kempf--Ness criterion \cite[Theorem 4, p.~1]{Losev}
the $SL_n$-orbit of $F$ is closed. The weights $a_i$ are linearly independent,
so a diagonal matrix maps $F$ to $F_n$, whose orbit is therefore closed.

(2) Half the Hessian matrix of $F_n$ is tridiagonal, with diagonal
$(x_2,x_3,\dots,x_n,3x_n)$ and off-diagonal entries $x_1,\dots,x_{n-1}$. Its
projective locus of rank at most two consists of
\begin{equation}\label{eq:hessianlocus}
 [e_1]\quad\text{and}\quad[ae_{n-1}+be_n]\ \text{with}\ a=0\ \text{or}\ a^2=3b^2 ,
\end{equation}
three collinear points and a fourth point off their line. Indeed, a tridiagonal
block whose off-diagonal entries are all nonzero has kernel of dimension at
most one, and the zero off-diagonal entries cut the matrix into such blocks; so
the rank is at least the number of nonzero $x_1,\dots,x_{n-1}$, and at most two
of them are nonzero. If exactly two, $x_i,x_j$ with $i<j$: when $j\ge i+2$ the
block on $\{i,i+1\}$ has determinant $-x_i^2$ (its first diagonal entry
$x_{i+1}$ vanishes) and the block of $x_j$ adds rank at least one; when $j=i+1$
and $i\ge2$, the singleton block $\{i-1\}$ has diagonal $x_i\ne0$, separate
from a block of rank at least two; when $i=1$, the block on $\{1,2,3\}$ has
determinant $-x_2^3\ne0$ for $n\ge5$, and for $n=4$ either this holds or
$x_4\ne0$ and the singleton $\{4\}$ adds one. If exactly one, $x_j$: for $2\le
j\le n-2$ the block $\{j,j+1\}$ has determinant $-x_j^2$ and the singleton
$\{j-1\}$ has diagonal $x_j$; for $j=1$ the block $\{1,2\}$ has rank two and
$x_n\ne0$ would add the last two diagonal entries, leaving $[e_1]$; for
$j=n-1$, with $a=x_{n-1}$, $b=x_n$, the singleton $\{n-2\}$ has entry $a$ and
the block $\left(\begin{smallmatrix}b&a\\a&3b\end{smallmatrix}\right)$ has rank
one exactly when $a^2=3b^2$. If none, the point is $[e_n]$, of rank two.

A linear transformation $T$ with $F_n\circ T=cF_n$ transforms the Hessian by
congruence, so it preserves \eqref{eq:hessianlocus}, hence fixes $[e_1]$, the
point off the line of the other three. In coordinate forms $\ell_i=(Tx)_i$
this says $\ell_1=ax_1+L(x_2,\dots,x_n)$ with $a\ne0$, while $\ell_2,\dots,
\ell_n$ do not involve $x_1$. The coefficient of $x_1^2$ in
$F_n(\ell)=cF_n(x)$ gives $a^2\ell_2=cx_2$; that of $x_1$ gives $2aL\ell_2=0$,
so $L=0$. Setting $x_1=0$, $(\ell_2,\dots,\ell_n)$ preserves up to a scalar the
chain cubic in $x_2,\dots,x_n$, with $\ell_2$ a multiple of $x_2$. If $n-1\ge4$
we repeat the argument; it remains to treat the ternary cubic
$G=y^2z+z^2w+w^3$ with $\ell=(by,uy+A,vy+B)$, $A,B$ independent linear forms
in $z,w$. The coefficient of $y$ in $G(\ell)$ is $vA^2+2uAB+3vB^2$; it
vanishes, so $u=v=0$. The coefficient of $y^2$ gives $b^2A=cz$; writing
$B=\alpha z+\beta w$ with $\beta\ne0$, the coefficient of $zw^2$ is
$3\alpha\beta^2=0$, so $\alpha=0$. Hence every such $T$ is diagonal. Dividing $T$ by a cube root of $c$, we may
take $c=1$. A diagonal $\diag(a_1,\dots,a_n)$ preserves $F_n$ exactly when $a_i^2a_{i+1}=1$
and $a_n^3=1$, that is $a_i=a_1^{(-2)^{i-1}}$ and $a_1^{3M}=1$; its
determinant is $a_1^S$ with $S=\sum_{i=0}^{n-1}(-2)^i=(1-(-2)^n)/3$, which is
odd and congruent to $n$ modulo three. So $\gcd(S,3M)=\nu$, and the kernel of
$\det$ has order $\nu$; when $\nu=3$ its elements are the scalars
$\omega I$, $\omega^3=1$, as $(-2)^i\equiv1\pmod3$, and these act trivially on
cubic forms.

(3) Since each fibre of the quotient map contains one closed orbit, a grading
element fixes $p_n$ exactly when it is $(\det a)^{-\gamma}$ for some $a$ in the
stabilizer of $F_n$. With $a_1$ of order $3M$, $(\det a)^{-\gamma}=
a_1^{-\gamma S}$ has order $3M/\gcd(3M,\gamma S)$. As $S$ is odd,
$\gcd(3M,\gamma S)=\gcd(3,\gamma S)$, which is $\gcd(3,3S)=3$ if $\nu=1$ and
$\gcd(3,S)=3$ if $\nu=3$ (then $S\equiv n\equiv0\pmod3$). So the grading
stabilizer of $p_n$ is $\mu_M$.

(4) The group $SL_n/\mu_\nu$ acts on $\Sym^3V^*$ and has trivial stabilizer at
$F_n$, whose orbit is closed. By Luna's slice theorem \cite[Section III.1,
Th\'eor\`eme du slice \'etale and Remarque 1$^\circ$, pp.~97--98]{Luna}, $X$ is smooth at $p_n$ with tangent space
$\Sym^3V^*/\mathfrak{sl}_n\cdot F_n$. Let $A=\diag(\zeta,\zeta^{-2},\dots,
\zeta^{(-2)^{n-1}})$ with $\zeta$ a primitive $M$-th root of unity; it fixes
$F_n$, and acts on $X$ as the grading element $(\det A)^{-\gamma}=\zeta^{-\gamma
S}$, a primitive $M$-th root of unity, whose fixed scheme is $X^\zeta$. The
monomials fixed by $A$ are exactly the $n$ chain monomials, by the valuation
argument at the end of the proof of Proposition~\ref{prop:graphsupport}; and
$\dim\mathfrak{sl}_n^A=n-1$, as $A$ has simple spectrum. Taking invariants is
exact, so $\dim(T_{p_n}X)^A=1$, and the fixed scheme is smooth of dimension one
at $p_n$ (proof of Proposition~\ref{prop:fixedcomponentpole}).
\end{proof}

\begin{theorem}\label{thm:plethysmPH3}
Let $n\ge4$ and $M=2^{n-1}$.
\begin{enumerate}
\item The support of the fixed scheme $X^{\mu_M}$ is the closure $Y$ of the
grading orbit of $p_n$, and $X$ is smooth along $Y\setminus\{0\}$.
\item $H_n$ has a simple pole at every primitive $M$-th root of unity $\zeta$,
with
\begin{equation}\label{eq:plethysmresidue}
 \mathcal C_n(\zeta):=\lim_{t\to\zeta}(1-t^M)H_n(t)
 =\frac1{\det\bigl(1-\zeta\mid N^*_{Y/X,\eta_Y}\bigr)}\ne0 .
\end{equation}
\item If $\eta$ is a root of unity whose order $q$ has two-part at least $M$,
then $H_n$ has a pole at $\eta$ if and only if $q=M$.
\item In every expression $H_n(t)=P(t)/\prod_j(1-t^{a_j})^{b_j}$ with
$P\in\C[t]$ and $a_j,b_j>0$, some $a_j$ is divisible by $M$.
\end{enumerate}
\end{theorem}

\begin{proof}
Let $p\ne0$ be a closed point fixed by a grading element of order $2^a\ge M$,
and $F$ a cubic in its closed orbit. Lemma~\ref{lem:lift} (with the odd
exponent $-\gamma$) gives a diagonal $A$ of two-power order fixing $F$ with
$\ord(\det A)=2^a$; Lemma~\ref{lem:span} gives the span condition; and
Theorem~\ref{thm:fullcubicsupport} gives $2^a\le M$, so $a=n-1$, and $F$ is
equivalent to $F_n$ under a permutation and a diagonal substitution. So
$p\in GL_n\cdot p_n$, which is the grading orbit of $p_n$. Conversely that
orbit is fixed by $\mu_M$ (Proposition~\ref{prop:chaincubic}(3)). The closure
$Y$ adds only the origin, so the support of $X^{\mu_M}$ is $Y$, a curve, and
$X$ is smooth along $Y\setminus\{0\}$ by Proposition~\ref{prop:chaincubic}(4)
and translation. Proposition~\ref{prop:fixedcomponentpole} with $\delta=1$ and
Lemma~\ref{lem:orbitclosure} ($e_Y=1$) give (2). For (3): if $q=2^ab$, $b$ odd,
$a\ge n-1$, and $q\ne M$, a nonzero point fixed by $\eta$ is fixed by its
two-primary part, of order $2^a\ge M$, hence lies on the grading orbit of
$p_n$, whose grading stabilizer is $\mu_M\not\ni\eta$. So $X^\eta$ is supported
at the origin, and Lemma~\ref{lem:isolated} applies. (4) follows from (2).
\end{proof}

\begin{corollary}\label{cor:plethysmPH3}
For $n\ge4$, every positive rational form of $H_n$ has modular index at least
$2^{n-1}$, while $\dim GL_m+\operatorname{ht}(\lambda)+\operatorname{ht}(\pi)
=m^2+n+1=O(n^6)$ and $\rank(x)=2n+m+1$. The existence clause of
\cite[Hypothesis 1.6.7]{GCT6}, and so also its specific clause on the minimal
positive form, fails for plethysm coefficients; so does
\cite[Hypothesis 3.3.4]{GCT6} for the plethysm problem.
\end{corollary}

